\subsection{Mutating External Scopes}

Reading an outer name follows LEGB automatically. Assigning to an outer name requires \texttt{global} (module namespace) or \texttt{nonlocal} (enclosing function cell). Without these declarations, assignment creates a local binding and shadows the outer name.

\subsubsection{Local Read Access Boundaries vs. the Shadowing Consequence of Assignment}

%<<<PROGRAM:programs/mut03>>>


Assignment inside the function creates a local \texttt{answer}; it does not mutate the module binding.

\subsubsection{Overriding Module Scope: The \texttt{global} Declaration Syntax and Module Namespace Mutation}

%<<<PROGRAM:programs/mut02>>>


\texttt{global} redirects \texttt{STORE\_GLOBAL} to the module namespace dictionary.

\subsubsection{Overriding Nested Intermediary Scopes: The \texttt{nonlocal} Declaration Syntax and Cell Reference Mutation}

%<<<PROGRAM:programs/mut01>>>


\texttt{nonlocal} writes through the closure cell, not a new local slot. Integer rebinding creates a new \texttt{int} object and updates the cell reference.

\subsubsection{\texttt{locals()} Caveats: Why the Displayed Local Namespace Is Not Always a Directly Writable Control Surface}

\texttt{locals()} returns a dict-like view of the current frame's local bindings. In optimized functions, mutating the returned dict may not affect fast-local slots. Treat \texttt{locals()} as an inspection aid, not a reliable control surface for rebinding.
